\section{Kubernetes (K8)}
\begin{mdframed}
	\textbf{K8 Objects:} Persistent entities which signal an intent (e.g. for something to be created on the cluster) \\
	\textbf{K8 Controller:} Tracks an Object and is responsible for bringing the current State closer to the desired State.
\end{mdframed}

\textbf{Pod:} Represents a process running on your cluster. Should contain one container, multiple are possible. \\
\textbf{Sidecars:} Sidecars are containers that run along the primary container in the same pod. Example use cases might be logging, security, data synchronization. \\
\textbf{Init Containers:} Similar to sidecars, but run and finish before app containers.
\textbf{Volume:} Assigns physical Storage to a Pod \\
\textbf{ReplicaSet:} Makes sure that a specified number of replica pods are running. In practice, deployments are used. \\
\textbf{Deployment:} Allows to manage one or multiple Pods. \\
\textbf{Service:} API Resource to expose logical set of Pods in the namespace. Acts as a load balancer (round-robin). \\
\textbf{Ingress:} Provides external Access to a Service. \\
\textbf{DaemonSet:} Ensures that all (or some) Nodes run a copy of a pod. This can be used for running supporting applications (logs, storage)
\textbf{ConfigMap:} Stores configuration data as key-value pairs. Used to decouple config from code. Mounted into pods as env vars or files. Not meant for sensitive data (use Secrets instead).
\textbf{Secret:} Like a ConfigMap but for sensitive data (passwords, tokens, certs). Values are only base64-encoded, \textit{not} encrypted by default (enable encryption at rest / use an external secret store). Mounted as env vars or files.
\begin{lstlisting}[language=yaml]
apiVersion: apps/v1
kind: Deployment
metadata:
  name: app
spec:
replicas: 3 # automatically deploys replicaset
selector:
  matchLabels:
    app: web
template:
  metadata: # pod
    labels:
      app: web
  spec:
    strategy:
      type: RollingUpdate # rolling update
      rollingUpdate:
        maxUnavailable: 1
        maxSurge: 1
    initContainers:
    - name: init
      image: busybox
      command: ['sh', '-c', 'sleep 10']
    containers: # container
    - name: web
      image: nginx
      ports:
      - containerPort: 80
    - name: sidecar
      image: busybox
      command: ['sh', '-c', 'while true; do sleep 30; done']
\end{lstlisting}

\subsection{Services}
A Service gives a stable virtual IP + DNS name to a dynamic set of Pods (selected by labels) and load-balances across them (round-robin). Types: \\
\textbf{ClusterIP:} (default) internal-only IP, reachable inside the cluster. \\
\textbf{NodePort:} opens a static port on every node, exposing the Service via \texttt{<NodeIP>:<port>}. \\
\textbf{LoadBalancer:} provisions an external (cloud) load balancer with a public IP that routes to the Service. \\
\textbf{Ingress} routes external HTTP(S) by host/path to Services (single entrypoint, needs an ingress controller). The newer \textbf{Gateway API} (\texttt{Gateway} + \texttt{HTTPRoute}) replaces Ingress with a more expressive, role-oriented model. \\
\textit{Port trap:} Service \texttt{port} (its own) vs \texttt{targetPort} (the container) vs \texttt{nodePort}; an Ingress/HTTPRoute backend must reference the \textit{Service} port, not the container port.

\subsection{Using ConfigMaps \& Secrets}
Inject into Pods as a single env var (\texttt{configMapKeyRef}/\texttt{secretKeyRef}), all keys at once (\texttt{envFrom}), or mounted as files (volume).
\begin{lstlisting}[language=yaml]
env:
- name: APIKEY
  valueFrom:
    configMapKeyRef: { name: cfg, key: api-key }
envFrom:
- configMapRef: { name: cfg }   # import every key as env var
volumeMounts:
- { name: v, mountPath: /etc/cfg }
volumes:
- name: v
  configMap: { name: cfg }      # or: secret: { secretName: s }
\end{lstlisting}

\subsection{LimitRange \& ResourceQuota}
\textbf{LimitRange:} min/max + default requests/limits \textit{per container} in a namespace; a Pod outside the range is \textit{rejected}, and a container omitting requests/limits gets the defaults auto-injected. \\
\textbf{ResourceQuota:} caps the \textit{summed} resources of a \textit{whole namespace}; a Pod that would exceed the total is rejected.

\subsection{Namespaces}
Used to separate resources. Only resources in same namespace can communicate directly.

\subsection{Rolling Updates}
Rolling updates can be used in order to ensure that enough pods are always running.
Rolling updates can be using maxUnavailable (maximum No. of Pods upgrading at the same time) and maxSurge (max No. of Pods allowed to run beyond specified No. of replica)

\subsection{Scheduling}
The \texttt{kube-scheduler} determines which nodes run which Pods. We can influence this decision process:

\begin{lstlisting}[language=yaml]
kind: Pod
spec:
  nodeSelector:
    disktype: ssd # this label needs to be in pod.spec
\end{lstlisting}

To evaluate if a Node is eligible to run a Pod, the following things are considered: Port availability, CPU \& Memory resources, available volumes, specified labels. Additionally, scoring is used to evaluate remaining nodes with criteria: pods of same service should be on different nodes, nodes with few used resources are prioritized, node affinity. \\
\textit{Taints} can also be applied on nodes and pods, pods won't be deployed on nodes with matching taints. \textit{Tolerations} can be used to make exceptions to taints. Types of taints: \textit{NoSchedule, PreferNoSchedule, NoExecute} \\
CLI: \texttt{kubectl taint nodes <n> gpu=true:NoSchedule} and \texttt{kubectl label nodes <n> disktype=ssd} (remove either by appending a trailing \texttt{-}). \textit{Gotcha:} a toleration only \textit{allows} a Pod onto a tainted node, it does not \textit{force} it there -- pair it with a \texttt{nodeSelector}/affinity to actually pin the Pod.
\begin{lstlisting}[language=yaml]
tolerations:
- { key: "gpu", operator: "Equal", value: "true", effect: "NoSchedule" }
\end{lstlisting}

\subsection{Probes \& Resources}
The kubelet uses probes (\texttt{httpGet}, \texttt{tcpSocket} or \texttt{exec}) to check container health: \\
\textbf{Liveness:} is it alive? On failure $\rightarrow$ restart the container. \\
\textbf{Readiness:} is it ready to serve? On failure $\rightarrow$ remove pod from Service endpoints (no traffic), but don't restart. \\
\textbf{Startup:} guards slow-starting apps; disables the other probes until it succeeds. Repeated liveness failures $\rightarrow$ container restarted with exponential back-off $\rightarrow$ \texttt{CrashLoopBackOff}. \\
\textbf{requests:} guaranteed minimum the scheduler reserves (used for placement). \textbf{limits:} hard ceiling -- CPU is throttled, memory over-limit $\rightarrow$ OOMKilled.
\begin{lstlisting}[language=yaml]
resources:
  requests: { cpu: "250m", memory: "64Mi" }
  limits:   { cpu: "500m", memory: "128Mi" }
livenessProbe:
  httpGet: { path: /healthz, port: 80 }
  initialDelaySeconds: 5
  periodSeconds: 10
readinessProbe:
  httpGet: { path: /ready, port: 80 }
\end{lstlisting}

\subsection{Commands}
Apply a manifest.yaml: \texttt{kubectl \{create|apply|replace\} -f manifest.yaml} \\
Generate YAML w/o applying: \texttt{kubectl create deploy web --image=nginx --dry-run=client -o yaml > d.yaml} \\
Imperative: \texttt{kubectl run}, \texttt{kubectl scale deploy x --replicas=4}, \texttt{kubectl expose} \\
Inspect: \texttt{kubectl get pods -A -o wide}, \texttt{describe node <n>}, \texttt{get nodes --show-labels} \\
Connecting to a Pod: \texttt{kubectl exec -it nginx-xxx -- sh} (drop \texttt{-it} for one-shot) \\
Rollout: \texttt{kubectl rollout \{status|history\} deploy x}, \texttt{rollout undo deploy x --to-revision=N} \\
Default namespace: \texttt{kubectl config set-context --current --namespace=lab}
